% Part: computability
% Chapter: computability-theory
% Section: ce-closed-cup-cap

\documentclass[../../../include/open-logic-section]{subfiles}

\begin{document}

\olfileid{cmp}{thy}{clo}

\olsection[संगणनक्षमपणे प्रगणनीय संचांचा संयोग आणि छेद]{संगणनक्षमपणे
  प्रगणनीय संच संयोग आणि छेद यांखाली बंद असतात}

पुढील प्रमेय संगणनक्षमपणे प्रगणनीय संचांच्या काही संवरण गुणधर्मांबद्दल आहे.

\begin{thm}
$A$ आणि $B$ हे संगणनक्षमपणे प्रगणनीय असतील, तर
$A \cap B$ आणि $A \cup B$ देखील तसेच असतात.
\end{thm}

\begin{proof}
\olref[eqc]{thm:ce-equiv} मुळे संगणनक्षमपणे प्रगणनीय संचांची
वेगवेगळी वैशिष्ट्ये वापरता येतात. त्याचे उदाहरण म्हणून काही
वेगवेगळ्या सिद्धता देऊ.

पहिल्या दोन सिद्धतांत दोन्ही संच रिकामे नाहीत असे धरा. एखादा
रिकामा असेल, तर त्यांचा छेद रिकामा आणि संयोग दुसरा संच असतो;
निष्कर्ष त्या वेळी थेट मिळतो.
%READERNOTE{OLCMP-019 — गोठवलेल्या इंग्रजीतील पहिल्या दोन सिद्धता दोन्ही c.e. संचांना सर्वत्र परिभाषित प्रगणक आहेत असे धरतात. c.e. च्या व्याख्येत रिकामा संचही आहे, पण तो अशा फलनाची व्याप्ती नसतो. रिकाम्या संचाची स्वतंत्र सोपी स्थिती येथे स्पष्ट केली आहे.}
पहिल्या सिद्धतेसाठी $A$ चे प्रगणन संगणनक्षम फलन~$f$ करते
आणि $B$ चे प्रगणन संगणनक्षम फलन~$g$ करते असे समजा.
पुढीलप्रमाणे घ्या:
\begin{align*}
h(x) & = \umin{y}{(f(y) = x \lor g(y) = x)} \text{ आणि}\\
j(x) & = \umin{y}{(f((y)_0) = x \land g((y)_1) = x)}.
\end{align*}
मग $A \cup B$ हे $h$ चे परिभाषाक्षेत्र आहे आणि
$A \cap B$ हे~$j$ चे परिभाषाक्षेत्र आहे.

\begin{explain}
संगणनाच्या भाषेत येथे पुढील गोष्ट घडते: $A$ आणि $B$ यांचे
प्रगणन करणाऱ्या प्रक्रिया दिल्या असता, घटक $x$ हा $A \cup B$
मध्ये आहे का हे एका तरी प्रगणनात~$x$ शोधून अर्धनिर्णीत करता
येते. तसेच $x$ हा $A \cap B$ मध्ये आहे का हे दोन्ही
प्रगणनांत एकाच वेळी~$x$ शोधून अर्धनिर्णीत करता येते.
\end{explain}

दुसऱ्या सिद्धतेसाठी पुन्हा $A$ चे प्रगणन~$f$ करते आणि
$B$ चे प्रगणन~$g$ करते असे समजा. पुढीलप्रमाणे घ्या:
\[
k(x) = \begin{cases}
f(x/2) & \text{जर $x$ सम असेल तर} \\
g((x-1)/2) & \text{जर $x$ विषम असेल तर.}
\end{cases}
\]
मग $k$ हे $A \cup B$ चे प्रगणन करते. कल्पना एवढीच की
$k$ हे $f$ आणि~$g$ यांच्या प्रगणनांचा आलटूनपालटून वापर करते.
$A \cap
B$ चे प्रगणन अधिक अवघड आहे. $A \cap B$ रिकामा असेल,
तर तो अर्थातच संगणनक्षमपणे प्रगणनीय आहे. अन्यथा $c$ हा
$A \cap B$ मधील कोणताही घटक घ्या आणि $l$ ची व्याख्या
पुढीलप्रमाणे करा:
\[
l(x) = \begin{cases}
f((x)_0) & \text{जर $f((x)_0) = g((x)_1)$ असेल तर} \\
c & \text{अन्यथा.}
\end{cases}
\]
संगणनाच्या भाषेत, $l$ हे $f$ आणि $g$ यांच्या प्रगणनांतील
घटकांच्या जोड्यांमधून जाते आणि सापडलेली प्रत्येक जुळणी परत
करते; अन्यथा ते~$c$ परत करून प्रगणनात प्रगती करत नाही.

शेवटच्या सिद्धतेसाठी $A$ हे आंशिक फलन $m(x)$ चे
\emph{परिभाषाक्षेत्र} आणि $B$ हे आंशिक फलन $n(x)$ चे
परिभाषाक्षेत्र आहे असे समजा. मग $A \cap B$ हे आंशिक फलन
$m(x) +
n(x)$ चे परिभाषाक्षेत्र आहे.

\begin{explain}
संगणनाच्या भाषेत, $A$ हा ज्यांवर $m$ थांबते त्या
मूल्यांचा संच असेल आणि $B$ हा ज्यांवर $n$ थांबते त्या
मूल्यांचा संच असेल, तर $A \cap B$ हा दोन्ही प्रक्रिया
थांबतात अशा मूल्यांचा संच आहे.
\end{explain}

$A \cup B$ हा थांबणाऱ्या आदानांचा संच म्हणून व्यक्त करणे
अधिक कठीण आहे, कारण $m$ आणि $n$ यांचे समांतर अनुकरण
करावे लागते. $d$ हा $m$ चा निर्देशांक आणि $e$ हा $n$
चा निर्देशांक असू द्या; दुसऱ्या शब्दांत, $m =
\cfind{d}$ आणि $n = \cfind{e}$. मग $A \cup B$ हे
पुढील फलनाचे परिभाषाक्षेत्र आहे:
\[
p(x) = \umin{y}{(T(d,x,y) \lor T(e,x,y))}.
\]
\begin{explain}
संगणनाच्या भाषेत, आदान $x$ वरील $p$ हे $m$ किंवा
$n$ यांपैकी एखाद्याचे थांबणारे संगणन शोधते आणि त्यांपैकी
एक सापडल्यास थांबते.
\end{explain}
\end{proof}

\end{document}
